\clearpage
\VerdierSemanticAnchor{front-publications}
\VerdierPhysicalPageAnchor{0010}
\chapter*{J.-L. Verdier 著作目録}
\addcontentsline{toc}{chapter}{J.-L. Verdier 著作目録}

\small
\begin{description}
\VerdierSemanticAnchor{front-publications-01}
\item[{[01]}] \emph{保型形式に付随する積分について}（G. Shimura による）、
Séminaire Bourbaki 1960--61, no. 216, pp. 216-1--216-27,
Benjamin (1966).

\VerdierSemanticAnchor{front-publications-02}
\item[{[02]}] \emph{Poincaré 双対定理}, C.R.A.S. vol. 256,
pp. 2084--2086, Gauthier-Villars (1963).

\VerdierSemanticAnchor{front-publications-03}
\item[{[03]}] \emph{数体のエタール・コホモロジーに関するセミナー}
（M. Artin との共著）, Woods Hole Conference on Algebraic Geometry,
A.M.S. (1964).

\VerdierSemanticAnchor{front-publications-04}
\item[{[04]}] \emph{射影系の本質的同値}, C.R.A.S. vol. 261,
pp. 4950--4953, Gauthier-Villars (1965).

\VerdierSemanticAnchor{front-publications-05}
\item[{[05]}] \emph{局所コンパクト空間の次元}, C.R.A.S. vol. 261,
pp. 5293--5296, Gauthier-Villars (1965).

\VerdierSemanticAnchor{front-publications-06}
\item[{[06]}] \emph{副有限群のコホモロジーにおける双対性}, Jean-Pierre
Serre, ``Galois Cohomology'' の付録, Lecture Notes in Mathematics 5,
pp. 183--206, Springer (1964), 第4版 (1973).

\VerdierSemanticAnchor{front-publications-07}
\item[{[07]}] \emph{局所コンパクト空間上の構成可能層}, C.R.A.S. vol. 262,
pp. 12--15, Gauthier-Villars (1966).

\VerdierSemanticAnchor{front-publications-08}
\item[{[08]}] \emph{局所コンパクト空間のコホモロジーにおける双対性},
Séminaire Bourbaki 1965--66, no. 300, pp. 300-01--300-13,
Benjamin (1966).

\VerdierSemanticAnchor{front-publications-09}
\item[{[09]}] \emph{エタール・コホモロジーにおける Lefschetz 不動点公式},
Conference on Local Fields, 1966年 Driebergen の NUFFIC Summer School,
pp. 199--214, Springer (1967).

\VerdierSemanticAnchor{front-publications-10}
\item[{[10]}] \emph{スキームのエタール・コホモロジーにおける双対定理},
Conference on Local Fields, 1966年 Driebergen の NUFFIC Summer School,
pp. 184--198, Springer (1967).

\VerdierSemanticAnchor{front-publications-11}
\item[{[11]}] \emph{連接層のねじれ逆像に対する基底変換}, Bombay
Colloquium on Algebraic Geometry in 1968, pp. 393--408, Oxford
University Press (1969).
\end{description}
\normalsize
